\subsubsection{Relativistic Quantum-Mechanical Setting}
The wave-equation section uses one fixed operator language throughout. With the differential conventions from Section \ref{sec:differential},
\begin{align*}
\partial&=\partial_t+\sigv\cdot\gradv,\\
\partial^*&=\partial_t-\sigv\cdot\gradv,\\
(i\partial)^*&=-i\partial^*.
\end{align*}
Recall also that the d'Alembertian is
\begin{align*}
\Box
&:=\partial_t^2-\gradv^2\notag\\
&=\partial\partial^*=\partial^*\partial\notag\\
&=\langle\!\langle\partial\rangle\!\rangle^2.
\end{align*}
All differential operators act left to right.

\subsubsection{What a Complex-Paravector State Represents}
Throughout the spinful part of the wave-equation section, the state is not a space-time event but an algebra-valued amplitude field:
\begin{equation*}
\psi:\R\times\R^3\to\C\oplus\C^3,
\qquad
\psi(t,\mathbf r)=\psi_s(t,\mathbf r)+\psi_{\mathbf v}(t,\mathbf r)\cdot\sigv.
\end{equation*}
The argument \((t,\mathbf r)\) tells where the state is evaluated in space-time, while the value \(\psi(t,\mathbf r)\) lives in the internal complex-paravector algebra. This distinction matters. The real paravector
\begin{equation*}
X=t+\mathbf r\cdot\sigv
\end{equation*}
describes an event or interval variable, whereas \(\psi\) is a quantum state amplitude defined \emph{over} such events.

At each fixed point \((t,\mathbf r)\), the field carries one complex scalar channel \(\psi_s\) and one complex vector channel \(\psi_{\mathbf v}\). Equivalently, it carries four complex local amplitudes, i.e. eight real local degrees of freedom. This is why the complex-paravector form can play the role usually assigned to a spinor carrier: it stores the same amount of local data, but arranges it as a scalar-plus-vector object adapted to the sigma algebra rather than as a column with separately named components.

The scalar and vector pieces should not be interpreted as separately measurable particles or classical subfields. They are amplitude channels that mix under the wave operators. The scalar part is the sigma-number-zero channel, while the vector part is the sigma-number-one channel that couples directly to \(\sigv\cdot\gradv\), to spin projectors, and to the first-order Dirac factorization. In that sense the vector part of \(\psi\) is not a literal velocity field; it is the part of the quantum amplitude that carries directional and spin-sensitive information.

The quantum-mechanical interpretation is therefore bilinear rather than componentwise. The raw field \(\psi\) is not itself an observable. Observable densities are extracted from sandwiches such as
\begin{equation*}
\rev{\psi}\psi
=\rho_\psi+\mathbf J_\psi\cdot\sigv,
\qquad
\rev{\psi}O\psi.
\end{equation*}
The scalar part gives the ordinary probability or expectation density, while the vector part records the associated flow, transport, or spin-orientation information. Thus the extra vector channels are not discarded because they are meaningless; they are discarded only when one asks for a scalar expectation value rather than for a current-like or polarization-like quantity.

This viewpoint also explains the later division of labor inside the section. The spin-0 Klein--Gordon subsections use \(\psi\in\C\) because a scalar carrier is sufficient there. Once one passes to spin-\(\tfrac12\), axis projectors, or first-order Dirac factorization, the complex-paravector carrier becomes the natural one because the sigma algebra must act nontrivially on the state itself. Later, when the manuscript writes a two-component field
\begin{equation*}
\bPsi=
\begin{Bmatrix}
\bPsi_1\\[2pt]
\bPsi_2
\end{Bmatrix},
\end{equation*}
it is packaging algebra-valued amplitudes into the auxiliary \(W\)-map factorization of the Dirac system; that notation is a structural reorganization of the same state content, not a redefinition of what the state represents.

\subsubsection{Complex-Paravector Bilinears and Bra-Ket Extraction}
To make the bra-ket notation precise in the paravector setting, write
\begin{equation*}
\begin{aligned}
\psi&=\psi_s+\psi_{\mathbf v}\cdot\sigv,
&
O&=O_s+O_{\mathbf v}\cdot\sigv,\\
Z&=Z_s+Z_{\mathbf v}\cdot\sigv,
&
X&=X_s+X_{\mathbf v}\cdot\sigv,
\end{aligned}
\end{equation*}
with \(\psi_s,Z_s\in\C\), \(\psi_{\mathbf v},Z_{\mathbf v}\in\C^3\), \(X_s\in\R\), \(X_{\mathbf v}\in\R^3\), and with operator-valued components acting left to right. The three involutions on the state are
\begin{equation*}
\begin{aligned}
\conj{\psi}&=\psi_s^*-\psi_{\mathbf v}^*\cdot\sigv,\\
\rev{\psi}&=\psi_s^*+\psi_{\mathbf v}^*\cdot\sigv,\\
\crev{\psi}&=\psi_s-\psi_{\mathbf v}\cdot\sigv.
\end{aligned}
\end{equation*}
The basic component products are collected below.
% These tables are algebraic component identities.  Their vector rows are not
% identified with probability currents without an accompanying field equation.
\newenvironment{bilineartable}[2][]{%
  \begin{table*}[p]
  \centering
  \caption{Components of $#2$.}%
  \if\relax\detokenize{#1}\relax\else\label{#1}\fi
  \renewcommand{\arraystretch}{1.10}%
  \setlength{\tabcolsep}{2pt}%
  \small
  \begin{tabular}{@{}r@{:\quad}>{\raggedright\arraybackslash}p{0.86\textwidth}@{}}
  \toprule
}{%
  \bottomrule
  \end{tabular}
  \end{table*}
}
\newcommand{\bcomponentrow}[2]{\eqtext{#1} & \(\displaystyle #2\) \\}

\noindent The direct action of a complex paravector operator on a complex paravector state is recorded in Table \ref{tab:components-O-psi}.
\begin{bilineartable}[tab:components-O-psi]{\pv{\mathcal{O}}\pv{\psi}}
\bcomponentrow{scalar}{\eqtext{scalar}(\pv{\mathcal{O}})\eqtext{scalar}(\pv{\psi})+\eqtext{vector}(\pv{\mathcal{O}})\cdot\eqtext{vector}(\pv{\psi})}
\bcomponentrow{vector}{\eqtext{scalar}(\pv{\mathcal{O}})\eqtext{vector}(\pv{\psi})+\eqtext{vector}(\pv{\mathcal{O}})\eqtext{scalar}(\pv{\psi})+i\bigl(\eqtext{vector}(\pv{\mathcal{O}})\times\eqtext{vector}(\pv{\psi})\bigr)}
\end{bilineartable}

\noindent The three basic ordered two-sided products with a generic complex
paravector
\(\pv{Z}=\eqtext{scalar}(\pv{Z})+\eqtext{vector}(\pv{Z})\cdot\sigv\) are given in Tables
\ref{tab:components-psistar-Z-psi}--\ref{tab:components-psistarR-Z-psi}.
\begin{bilineartable}[tab:components-psistar-Z-psi]{\conj{\pv{\psi}}\pv{Z}\pv{\psi}}
\bcomponentrow{scalar}{\begin{aligned}[t]
&\eqtext{scalar}(\pv{Z})\bigl(|\eqtext{scalar}(\pv{\psi})|^2-\eqtext{vector}(\pv{\psi})^*\cdot\eqtext{vector}(\pv{\psi})\bigr)\\
&\quad+\eqtext{scalar}(\pv{\psi})^*(\eqtext{vector}(\pv{Z})\cdot\eqtext{vector}(\pv{\psi}))
-\eqtext{scalar}(\pv{\psi})(\eqtext{vector}(\pv{Z})\cdot\eqtext{vector}(\pv{\psi})^*)\\
&\quad-i\,\eqtext{vector}(\pv{\psi})^*\cdot(\eqtext{vector}(\pv{Z})\times\eqtext{vector}(\pv{\psi}))
\end{aligned}}
\bcomponentrow{vector}{\begin{aligned}[t]
&\eqtext{scalar}(\pv{Z})\bigl(\eqtext{scalar}(\pv{\psi})^*\eqtext{vector}(\pv{\psi})-\eqtext{scalar}(\pv{\psi})\eqtext{vector}(\pv{\psi})^*
-i\,\eqtext{vector}(\pv{\psi})^*\times\eqtext{vector}(\pv{\psi})\bigr)\\
&\quad+\bigl(|\eqtext{scalar}(\pv{\psi})|^2+\eqtext{vector}(\pv{\psi})^*\cdot\eqtext{vector}(\pv{\psi})\bigr)\eqtext{vector}(\pv{Z})\\
&\quad-(\eqtext{vector}(\pv{Z})\cdot\eqtext{vector}(\pv{\psi})^*)\eqtext{vector}(\pv{\psi})
-(\eqtext{vector}(\pv{Z})\cdot\eqtext{vector}(\pv{\psi}))\eqtext{vector}(\pv{\psi})^*\\
&\quad+i\bigl(\eqtext{scalar}(\pv{\psi})^* \eqtext{vector}(\pv{Z})\times\eqtext{vector}(\pv{\psi})
+\eqtext{scalar}(\pv{\psi}) \eqtext{vector}(\pv{Z})\times\eqtext{vector}(\pv{\psi})^*\bigr)
\end{aligned}}
\end{bilineartable}

\begin{bilineartable}[tab:components-psiR-Z-psi]{\herm{\pv{\psi}}\pv{Z}\pv{\psi}}
\bcomponentrow{scalar}{\begin{aligned}[t]
&\eqtext{scalar}(\pv{Z})\bigl(|\eqtext{scalar}(\pv{\psi})|^2+\eqtext{vector}(\pv{\psi})^*\cdot\eqtext{vector}(\pv{\psi})\bigr)\\
&\quad+\eqtext{scalar}(\pv{\psi})^*(\eqtext{vector}(\pv{Z})\cdot\eqtext{vector}(\pv{\psi}))
+\eqtext{scalar}(\pv{\psi})(\eqtext{vector}(\pv{Z})\cdot\eqtext{vector}(\pv{\psi})^*)\\
&\quad+i\,\eqtext{vector}(\pv{\psi})^*\cdot(\eqtext{vector}(\pv{Z})\times\eqtext{vector}(\pv{\psi}))
\end{aligned}}
\bcomponentrow{vector}{\begin{aligned}[t]
&\eqtext{scalar}(\pv{Z})\bigl(\eqtext{scalar}(\pv{\psi})^*\eqtext{vector}(\pv{\psi})+\eqtext{scalar}(\pv{\psi})\eqtext{vector}(\pv{\psi})^*
+i\,\eqtext{vector}(\pv{\psi})^*\times\eqtext{vector}(\pv{\psi})\bigr)\\
&\quad+\bigl(|\eqtext{scalar}(\pv{\psi})|^2-\eqtext{vector}(\pv{\psi})^*\cdot\eqtext{vector}(\pv{\psi})\bigr)\eqtext{vector}(\pv{Z})\\
&\quad+(\eqtext{vector}(\pv{Z})\cdot\eqtext{vector}(\pv{\psi})^*)\eqtext{vector}(\pv{\psi})
+(\eqtext{vector}(\pv{Z})\cdot\eqtext{vector}(\pv{\psi}))\eqtext{vector}(\pv{\psi})^*\\
&\quad+i\bigl(\eqtext{scalar}(\pv{\psi})^* \eqtext{vector}(\pv{Z})\times\eqtext{vector}(\pv{\psi})
-\eqtext{scalar}(\pv{\psi}) \eqtext{vector}(\pv{Z})\times\eqtext{vector}(\pv{\psi})^*\bigr)
\end{aligned}}
\end{bilineartable}

\begin{bilineartable}[tab:components-psistarR-Z-psi]{\cherm{\pv{\psi}}\pv{Z}\pv{\psi}}
\bcomponentrow{scalar}{\eqtext{scalar}(\pv{Z})\bigl(\eqtext{scalar}(\pv{\psi})^2-\eqtext{vector}(\pv{\psi})\cdot\eqtext{vector}(\pv{\psi})\bigr)}
\bcomponentrow{vector}{\begin{aligned}[t]
&\bigl(\eqtext{scalar}(\pv{\psi})^2+\eqtext{vector}(\pv{\psi})\cdot\eqtext{vector}(\pv{\psi})\bigr)\eqtext{vector}(\pv{Z})\\
&\quad-2(\eqtext{vector}(\pv{Z})\cdot\eqtext{vector}(\pv{\psi}))\eqtext{vector}(\pv{\psi})
+2i\,\eqtext{scalar}(\pv{\psi}) \eqtext{vector}(\pv{Z})\times\eqtext{vector}(\pv{\psi})
\end{aligned}}
\end{bilineartable}

\noindent Tables \ref{tab:components-psistar-X-psi}--\ref{tab:components-psistarR-X-psi} are obtained by specialization to a real paravector \(\pv{X}=\eqtext{scalar}(\pv{X})+\eqtext{vector}(\pv{X})\cdot\sigv\), with \(\eqtext{scalar}(\pv{X})\in\R\) and \(\eqtext{vector}(\pv{X})\in\R^3\).
\begin{bilineartable}[tab:components-psistar-X-psi]{\conj{\pv{\psi}}\pv{X}\pv{\psi}}
\bcomponentrow{scalar}{\begin{aligned}[t]
&\eqtext{scalar}(\pv{X})\bigl(|\eqtext{scalar}(\pv{\psi})|^2-\eqtext{vector}(\pv{\psi})^*\cdot\eqtext{vector}(\pv{\psi})\bigr)\\
&\quad+\eqtext{scalar}(\pv{\psi})^*(\eqtext{vector}(\pv{X})\cdot\eqtext{vector}(\pv{\psi}))
-\eqtext{scalar}(\pv{\psi})(\eqtext{vector}(\pv{X})\cdot\eqtext{vector}(\pv{\psi})^*)\\
&\quad-i\,\eqtext{vector}(\pv{\psi})^*\cdot(\eqtext{vector}(\pv{X})\times\eqtext{vector}(\pv{\psi}))
\end{aligned}}
\bcomponentrow{vector}{\begin{aligned}[t]
&\eqtext{scalar}(\pv{X})\bigl(\eqtext{scalar}(\pv{\psi})^*\eqtext{vector}(\pv{\psi})-\eqtext{scalar}(\pv{\psi})\eqtext{vector}(\pv{\psi})^*
-i\,\eqtext{vector}(\pv{\psi})^*\times\eqtext{vector}(\pv{\psi})\bigr)\\
&\quad+\bigl(|\eqtext{scalar}(\pv{\psi})|^2+\eqtext{vector}(\pv{\psi})^*\cdot\eqtext{vector}(\pv{\psi})\bigr)\eqtext{vector}(\pv{X})\\
&\quad-(\eqtext{vector}(\pv{X})\cdot\eqtext{vector}(\pv{\psi})^*)\eqtext{vector}(\pv{\psi})
-(\eqtext{vector}(\pv{X})\cdot\eqtext{vector}(\pv{\psi}))\eqtext{vector}(\pv{\psi})^*\\
&\quad+i\bigl(\eqtext{scalar}(\pv{\psi})^* \eqtext{vector}(\pv{X})\times\eqtext{vector}(\pv{\psi})
+\eqtext{scalar}(\pv{\psi}) \eqtext{vector}(\pv{X})\times\eqtext{vector}(\pv{\psi})^*\bigr)
\end{aligned}}
\end{bilineartable}

\begin{bilineartable}[tab:components-psiR-X-psi]{\herm{\pv{\psi}}\pv{X}\pv{\psi}}
\bcomponentrow{scalar}{\begin{aligned}[t]
&\eqtext{scalar}(\pv{X})\bigl(|\eqtext{scalar}(\pv{\psi})|^2+\eqtext{vector}(\pv{\psi})^*\cdot\eqtext{vector}(\pv{\psi})\bigr)\\
&\quad+\eqtext{scalar}(\pv{\psi})^*(\eqtext{vector}(\pv{X})\cdot\eqtext{vector}(\pv{\psi}))
+\eqtext{scalar}(\pv{\psi})(\eqtext{vector}(\pv{X})\cdot\eqtext{vector}(\pv{\psi})^*)\\
&\quad+i\,\eqtext{vector}(\pv{\psi})^*\cdot(\eqtext{vector}(\pv{X})\times\eqtext{vector}(\pv{\psi}))
\end{aligned}}
\bcomponentrow{vector}{\begin{aligned}[t]
&\eqtext{scalar}(\pv{X})\bigl(\eqtext{scalar}(\pv{\psi})^*\eqtext{vector}(\pv{\psi})+\eqtext{scalar}(\pv{\psi})\eqtext{vector}(\pv{\psi})^*
+i\,\eqtext{vector}(\pv{\psi})^*\times\eqtext{vector}(\pv{\psi})\bigr)\\
&\quad+\bigl(|\eqtext{scalar}(\pv{\psi})|^2-\eqtext{vector}(\pv{\psi})^*\cdot\eqtext{vector}(\pv{\psi})\bigr)\eqtext{vector}(\pv{X})\\
&\quad+(\eqtext{vector}(\pv{X})\cdot\eqtext{vector}(\pv{\psi})^*)\eqtext{vector}(\pv{\psi})
+(\eqtext{vector}(\pv{X})\cdot\eqtext{vector}(\pv{\psi}))\eqtext{vector}(\pv{\psi})^*\\
&\quad+i\bigl(\eqtext{scalar}(\pv{\psi})^* \eqtext{vector}(\pv{X})\times\eqtext{vector}(\pv{\psi})
-\eqtext{scalar}(\pv{\psi}) \eqtext{vector}(\pv{X})\times\eqtext{vector}(\pv{\psi})^*\bigr)
\end{aligned}}
\end{bilineartable}

\begin{bilineartable}[tab:components-psistarR-X-psi]{\cherm{\pv{\psi}}\pv{X}\pv{\psi}}
\bcomponentrow{scalar}{\eqtext{scalar}(\pv{X})\bigl(\eqtext{scalar}(\pv{\psi})^2-\eqtext{vector}(\pv{\psi})\cdot\eqtext{vector}(\pv{\psi})\bigr)}
\bcomponentrow{vector}{\begin{aligned}[t]
&\bigl(\eqtext{scalar}(\pv{\psi})^2+\eqtext{vector}(\pv{\psi})\cdot\eqtext{vector}(\pv{\psi})\bigr)\eqtext{vector}(\pv{X})\\
&\quad-2(\eqtext{vector}(\pv{X})\cdot\eqtext{vector}(\pv{\psi}))\eqtext{vector}(\pv{\psi})
+2i\,\eqtext{scalar}(\pv{\psi}) \eqtext{vector}(\pv{X})\times\eqtext{vector}(\pv{\psi})
\end{aligned}}
\end{bilineartable}

\noindent For a scalar function \(S\), the same three products are reduced to
Tables \ref{tab:components-psistar-S-psi}--\ref{tab:components-psistarR-S-psi}.
\begin{bilineartable}[tab:components-psistar-S-psi]{\conj{\pv{\psi}}S\pv{\psi}}
\bcomponentrow{scalar}{S\bigl(|\eqtext{scalar}(\pv{\psi})|^2-\eqtext{vector}(\pv{\psi})^*\cdot\eqtext{vector}(\pv{\psi})\bigr)}
\bcomponentrow{vector}{S\bigl(\eqtext{scalar}(\pv{\psi})^*\eqtext{vector}(\pv{\psi})-\eqtext{scalar}(\pv{\psi})\eqtext{vector}(\pv{\psi})^*-i\,\eqtext{vector}(\pv{\psi})^*\times\eqtext{vector}(\pv{\psi})\bigr)}
\end{bilineartable}

\begin{bilineartable}[tab:components-psiR-S-psi]{\herm{\pv{\psi}}S\pv{\psi}}
\bcomponentrow{scalar}{S\bigl(|\eqtext{scalar}(\pv{\psi})|^2+\eqtext{vector}(\pv{\psi})^*\cdot\eqtext{vector}(\pv{\psi})\bigr)}
\bcomponentrow{vector}{S\bigl(\eqtext{scalar}(\pv{\psi})^*\eqtext{vector}(\pv{\psi})+\eqtext{scalar}(\pv{\psi})\eqtext{vector}(\pv{\psi})^*+i\,\eqtext{vector}(\pv{\psi})^*\times\eqtext{vector}(\pv{\psi})\bigr)}
\end{bilineartable}

\begin{bilineartable}[tab:components-psistarR-S-psi]{\cherm{\pv{\psi}}S\pv{\psi}}
\bcomponentrow{scalar}{S\bigl(\eqtext{scalar}(\pv{\psi})^2-\eqtext{vector}(\pv{\psi})\cdot\eqtext{vector}(\pv{\psi})\bigr)}
\bcomponentrow{vector}{0}
\end{bilineartable}


The reverse sandwich with a real scalar is the key probabilistic object. For \(S\in\R\),
\begin{equation*}
\rev{\psi}S\psi
=S\bigl(\rho_\psi+\mathbf J_\psi\cdot\sigv\bigr),
\end{equation*}
where
\begin{equation*}
\rho_\psi:=|\psi_s|^2+\psi_{\mathbf v}^*\cdot\psi_{\mathbf v},
\qquad
\mathbf J_\psi:=\psi_s^*\psi_{\mathbf v}
+\psi_s\psi_{\mathbf v}^*
+i\,\psi_{\mathbf v}^*\times\psi_{\mathbf v}.
\end{equation*}
Thus \(\rev{\psi}\psi\) is naturally a density-current paravector: the scalar part is the positive density, while the vector part carries the associated directional information. More generally, the raw bilinear \(\rev{\psi}O\psi\) is a complex paravector, so before any extraction it contains one complex scalar and one complex vector, that is, four real component channels. The scalar channel gives the ordinary expectation value; the remaining channels can encode currents, fluxes, spin-orientation data, or other geometric bilinears depending on the operator.

At fixed time, the natural normalization is therefore
\begin{equation*}
\langle \psi\mid\psi\rangle
:=\int_{\R^3}\eqtext{scalar}(\rev{\psi}\psi)\,\dd^3r
=\int_{\R^3}\rho_\psi\,\dd^3r
=1.
\end{equation*}
For a complex-paravector operator \(O\), define
\begin{equation*}
\langle \psi\mid O\mid\psi\rangle
:=\int_{\R^3}\eqtext{scalar}(\rev{\psi}O\psi)\,\dd^3r.
\end{equation*}
If the state is not already normalized, the corresponding expectation map is
\begin{equation*}
\mathbb E_\psi[O]
:=
\frac{\int_{\R^3}\eqtext{scalar}(\rev{\psi}O\psi)\,\dd^3r}
{\int_{\R^3}\eqtext{scalar}(\rev{\psi}\psi)\,\dd^3r}.
\end{equation*}
If \(O\psi=\phi_s+\phi_{\mathbf v}\cdot\sigv\), then the extracted scalar channel is simply
\begin{equation*}
\eqtext{scalar}(\rev{\psi}O\psi)
=\psi_s^*\phi_s+\psi_{\mathbf v}^*\cdot\phi_{\mathbf v}.
\end{equation*}
When \(O\) is reverse-self-adjoint under the spatial integral and the boundary terms are benign, this scalar expectation is real.

By contrast,
\begin{equation*}
\crev{\psi}\psi
=\psi_s^2-\psi_{\mathbf v}\cdot\psi_{\mathbf v}
=\langle\!\langle\psi\rangle\!\rangle^2
\end{equation*}
is the STA metric bilinear. It is algebraically useful, but it is not positive-definite and so it should not be used as the probabilistic normalization. In this framework the single bra-ket bracket \(\langle\cdot\mid\cdot\rangle\) is therefore best reserved for the scalar expectation map built from reversion, while \(\langle\!\langle\cdot\rangle\!\rangle^2\) continues to denote the paravector metric.

\subsubsection{NEW CONTENT FOR REVIEW}
Standard physics does not assign spin from component count alone. The physically relevant data are
\begin{equation*}
(\text{transformation law},\ \text{field equation},\ \text{constraints}).
\end{equation*}
In standard Lorentz language,
\begin{align*}
\text{spin }0:&\quad (0,0),\\
\text{spin }\tfrac12:&\quad (\tfrac12,0)\oplus(0,\tfrac12),\\
\text{spin }1:&\quad (1,0)\oplus(0,1)\quad \text{for }F,\\
&\quad (\tfrac12,\tfrac12)\quad \text{for }A^\mu.
\end{align*}
So the algebraic carrier spaces
\begin{equation*}
\C,\qquad
\C\oplus\C^3,\qquad
\varnothing\oplus\C^3
\end{equation*}
do not determine the physical spin by themselves. They become spin sectors only after the operator law is fixed.

\noindent\textbf{Bilinear pairings.}
The full reverse product is the one that directly produces the density-current paravector:
\begin{equation*}
\rev{\psi}\psi
=\rho_\psi+\mathbf J_\psi\cdot\sigv,
\qquad
\rho_\psi
=|\psi_s|^2+\psi_{\mathbf v}^*\cdot\psi_{\mathbf v}.
\end{equation*}
For two fields \(\phi,\psi\in\C\oplus\C^3\), the three scalar pairings are
\begin{align*}
\eqtext{scalar}(\rev{\phi}\psi)
&=\phi_s^*\psi_s+\phi_{\mathbf v}^*\cdot\psi_{\mathbf v}
&& \text{Hermitian-type pairing},\\
\eqtext{scalar}(\crev{\phi}\psi)
&=\phi_s\psi_s-\phi_{\mathbf v}\cdot\psi_{\mathbf v}
&& \text{metric pairing},\\
\eqtext{scalar}(\conj{\phi}\psi)
&=\phi_s^*\psi_s-\phi_{\mathbf v}^*\cdot\psi_{\mathbf v}
&& \text{indefinite companion}.
\end{align*}
For \(\phi=\psi\), the metric pairing reduces to
\begin{equation*}
\eqtext{scalar}(\crev{\psi}\psi)
=\crev{\psi}\psi
=\langle\!\langle\psi\rangle\!\rangle^2.
\end{equation*}
Only the reverse pairing is being used here as the quantum bra-ket pairing. The metric pairing is the intrinsic STA quadratic form used in factorization identities, and the companion pairing is algebraically well-defined but does not yet carry an independent standard quantum interpretation in this manuscript.

Accordingly, for the spatially integrated bracket
\begin{equation}
\langle \phi \mid \psi \rangle
:=\int_{\R^3}\eqtext{scalar}(\rev{\phi}\psi)\,\dd^3r,
\qquad
\langle \phi \mid O\psi \rangle
=\langle O^\dagger\phi \mid \psi \rangle,
\label{eq:adjoint-paravector-braket}
\end{equation}
the operator adjoint \(O^\dagger\) is the reverse-adjoint determined by this pairing, up to the usual boundary terms when derivatives are present. For left multiplication by a fixed paravector \(Z\), the adjoint is left multiplication by \(\rev{Z}\), since \(\rev{(\rev{Z}\phi)}=\rev{\phi}Z\). The companion vector channel is also worth naming:
\begin{equation*}
\mathbf J_{O,\psi}
:=\eqtext{vector}(\rev{\psi}O\psi).
\end{equation*}
The case \(O=1\) gives the density-current pair \((\rho_\psi,\mathbf J_\psi)\), while more general \(O\) produce current-, flux-, or polarization-type vector bilinears alongside the ordinary scalar expectation value.

\noindent\textbf{Generalized operator products.}
The following identities are algebraic. They describe operator composition on a paravector-valued state, but by themselves they do not assign a physical spin.

For operator-valued components, the commutator no longer collapses to the pure cross-product law of Eq.\ \eqref{eq:paravector-commutator}. Write
\begin{equation*}
O=O_s+O_{\mathbf v}\cdot\sigv,
\qquad
O'=O_s'+O_{\mathbf v}'\cdot\sigv,
\end{equation*}
with component operators commuting with the fixed basis elements and acting left to right on the state. Then, with dot and cross products taken componentwise and in the written operator order,
\begin{align}
[O,O']\psi
&=\Bigl([O_s,O_s']
+O_{\mathbf v}\cdot O_{\mathbf v}'
-O_{\mathbf v}'\cdot O_{\mathbf v}\Bigr)\psi\notag\\
&\quad+\Bigl([O_s,O_{\mathbf v}']
+[O_{\mathbf v},O_s']\notag\\
&\qquad\qquad
+i\bigl(O_{\mathbf v}\times O_{\mathbf v}'
-O_{\mathbf v}'\times O_{\mathbf v}\bigr)\Bigr)\cdot\sigv\,\psi,
\label{eq:operator-commutator-action}
\end{align}
where the scalar-vector commutators are understood componentwise. Unlike the purely algebraic paravector commutator, Eq.\ \eqref{eq:operator-commutator-action} generally has both scalar and vector parts. Only when all component operators mutually commute does it reduce to the simpler law
\begin{equation*}
[O,O']\psi
=2i\bigl(O_{\mathbf v}\times O_{\mathbf v}'\bigr)\cdot\sigv\,\psi.
\end{equation*}

The matching anticommutator is
\begin{align}
\operatorname{acom}(O,O')\psi
&:=(OO'+O'O)\psi\notag\\
&=\Bigl(\operatorname{acom}(O_s,O_s')
+O_{\mathbf v}\cdot O_{\mathbf v}'
+O_{\mathbf v}'\cdot O_{\mathbf v}\Bigr)\psi\notag\\
&\quad+\Bigl(\operatorname{acom}(O_s,O_{\mathbf v}')\notag\\
&\qquad\qquad
+\operatorname{acom}(O_{\mathbf v},O_s')\notag\\
&\qquad\qquad
+i\bigl(O_{\mathbf v}\times O_{\mathbf v}'
+O_{\mathbf v}'\times O_{\mathbf v}\bigr)\Bigr)\cdot\sigv\,\psi,
\label{eq:operator-anticommutator-action}
\end{align}
so that
\begin{align*}
OO'
&=\tfrac12\bigl(\operatorname{acom}(O,O')+[O,O']\bigr).
\end{align*}
This is the natural companion to the operator commutator formula and is the form from which covariance and uncertainty identities would be built if that development is expanded later.

\noindent\textbf{Carrier closure.}
The same algebra shows which carrier spaces are actually closed under which operators:
\begin{align}
\partial:\C&\to\C\oplus\C^3,\notag\\
\partial:\varnothing\oplus\C^3&\to\C\oplus\C^3,\notag\\
\partial:\C\oplus\C^3&\to\C\oplus\C^3.
\label{eq:carrier-closure}
\end{align}
Indeed,
\begin{align*}
\partial\psi_s
&=(\partial_t\psi_s)+(\gradv\psi_s)\cdot\sigv,\\
\partial(\psi_{\mathbf v}\cdot\sigv)
&=(\gradv\cdot\psi_{\mathbf v})
+\bigl(\partial_t\psi_{\mathbf v}+i\,\gradv\times\psi_{\mathbf v}\bigr)\cdot\sigv.
\end{align*}
So the full complex-paravector carrier is the minimal space closed under first-order paravector operators such as \(\partial\) and \(E\pm\sigv\cdot\mathbf P\). The scalar subspace closes only after passing to scalar second-order combinations such as \(\Box+m^2\). A pure-vector specialization is algebraically meaningful, but it is not preserved by generic first-order paravector action and therefore requires Maxwell- or Proca-type constraints if it is to describe a spin-1 field. A real-paravector specialization \(\psi\in\R\oplus\R^3\) is also possible algebraically, but it is not a general quantum carrier because it is not closed under multiplication by the phase \(e^{i\alpha}\).

With those algebraic preliminaries fixed, the organizing principle is this: one first chooses a carrier field for \(\psi\), then restricts to the operators that preserve that carrier or define its field equation, and then lists the bilinear pairings that extract the physically relevant observables from that carrier. The bilinears need not take values in the same carrier field; in general they land in scalar, vector, or paravector observable channels.

With that understood, the physically consistent sector-by-sector summary is the following.

\noindent\textbf{1. Spin \(0\): complex scalar Klein--Gordon sector}
\begin{align*}
\text{carrier:}\quad
&\psi:\R^{1,3}\to\C,\\
\text{representation:}\quad
&(0,0),\\
\text{field equation:}\quad
&(\Box+m^2)\psi=0,\\
\text{operator class:}\quad
&O_s:\C\to\C,\\
\text{examples:}\quad
&\Box+m^2,\ \partial_t\pm i\Omega,\\
&-E^2+\mathbf P^2+m^2,\ldots,\\
\text{spin projector:}\quad
&\Pi(\pm)\psi\notin\C.
\end{align*}
The scalar pairings reduce to
\begin{align*}
\eqtext{scalar}(\rev{\phi}\psi)
&=\phi^*\psi,\\
\eqtext{scalar}(\crev{\phi}\psi)
&=\phi\psi,\\
\eqtext{scalar}(\conj{\phi}\psi)
&=\phi^*\psi.
\end{align*}
So the reverse and sigma-conjugate pairings coincide on the scalar sector, while the metric pairing is different. The ordinary complex sesquilinear product \(\phi^*\psi\) is the known scalar-field pairing, but the standard relativistic conserved Klein--Gordon current is derivative-built rather than simply \(\psi^*\psi\):
\begin{equation*}
j^\mu_{\mathrm{KG}}
:=\frac{i}{2m}\bigl(\psi^*\partial^\mu\psi-(\partial^\mu\psi^*)\psi\bigr).
\end{equation*}
Thus \(\phi^*\psi\) is the natural Hermitian scalar pairing, while \(\phi\psi\) is an algebraic metric pairing and not a probability density.

\noindent\textbf{2. Spin \(\tfrac12\): Dirac/Pauli sector}
\begin{align*}
\text{algebraic block:}\quad
&\psi:\R^{1,3}\to\C\oplus\C^3,\\
\text{physical field:}\quad
&\bPsi=
\begin{Bmatrix}
\bPsi_1\\[2pt]
\bPsi_2
\end{Bmatrix},
\qquad
\bPsi_1,\bPsi_2\in\C\oplus\C^3,\\
\text{representation:}\quad
&(\tfrac12,0)\oplus(0,\tfrac12),\\
\text{field equations:}\quad
&(\Wmat(i\partial)-m\I)\bPsi=0,\\
&(\Wmat(i\partial-q\Phi^*)-m\I)\bPsi=0,\\
\text{operator class:}\quad
&O=O_s+O_{\mathbf v}\cdot\sigv,\\
&\text{acting on each paravector block},\\
\text{examples:}\quad
&\partial,\ \partial^*,\\
&E\pm\sigv\cdot\mathbf P,\ \Pi(\pm),\\
\text{spin projector:}\quad
&\Pi(\pm):\C\oplus\C^3\to\C\oplus\C^3,\\
&\bPsi(\pm):=\Pi(\pm)\bPsi,\\
&(\mathbf u\cdot\sigv)\bPsi(\pm)=\pm\bPsi(\pm).
\end{align*}
The basic one-block pairings are
\begin{align*}
\eqtext{scalar}(\rev{\phi}\psi)
&=\phi_s^*\psi_s+\phi_{\mathbf v}^*\cdot\psi_{\mathbf v},\\
\mathbf J_{\phi,\psi}
&:=\eqtext{vector}(\rev{\phi}\psi),\\
\eqtext{scalar}(\crev{\phi}\psi)
&=\phi_s\psi_s-\phi_{\mathbf v}\cdot\psi_{\mathbf v},\\
\eqtext{scalar}(\conj{\phi}\psi)
&=\phi_s^*\psi_s-\phi_{\mathbf v}^*\cdot\psi_{\mathbf v}.
\end{align*}
For the two-block Dirac carrier, the natural extension is
\begin{equation*}
\langle \{\Xi_1;\Xi_2\} \mid \{\Psi_1;\Psi_2\} \rangle
:=\int_{\R^3}
\sum_{n=1}^{2}\eqtext{scalar}(\rev{\Xi_n}\Psi_n)\,\dd^3r.
\end{equation*}
Here the reverse pairing is the Hermitian-type pairing used for expectation values and for current or polarization extraction. The metric pairing is the one used in factorization identities such as \(Z\crev{Z}\) and \(\Wmat(Z)^2\). The companion pairing is indefinite and no unique standard physical observable has been attached to it here.

The key physics point is that the one-block complex paravector is the minimal algebraic block preserved by first-order paravector action, but its scalar-plus-vector component count does not by itself make it an irreducible spin-\(\tfrac12\) field in the standard Lorentz sense. In this manuscript the physical spin-\(\tfrac12\) assignment belongs to the later Dirac system built from those paravector blocks.

\noindent\textbf{3. Spin \(1\): vector or field-strength sector}
\begin{align*}
\text{carrier:}\quad
&F:\R^{1,3}\to\varnothing\oplus\C^3,\\
\text{main example:}\quad
&F=(\mathbf E+i\mathbf B)\cdot\sigv,\\
\text{representation:}\quad
&(1,0)\oplus(0,1)\quad\text{for }F,\\
&(\tfrac12,\tfrac12)\quad\text{for a vector potential }A^\mu,\\
\text{field equation:}\quad
&\partial F=D^*,\\
\text{operator class:}\quad
&\partial,\ \partial^*,\ \Box,\ \text{and Maxwell/Proca constraints},\\
\text{closure facts:}\quad
&\partial,\partial^*:\varnothing\oplus\C^3\to\C\oplus\C^3,\\
&\Box:\varnothing\oplus\C^3\to\varnothing\oplus\C^3,\\
\text{spin projector:}\quad
&\Pi(\pm)F\notin\varnothing\oplus\C^3.
\end{align*}
So the two-projector split \(\Pi(\pm)\) is not the standard spin-\(1\) decomposition. The standard polarization counts are
\begin{align*}
\text{massive spin }1:&\quad m_s=-1,0,1,\\
\text{massless spin }1:&\quad \lambda=\pm1.
\end{align*}
The bilinear quantities already present in the manuscript are
\begin{align*}
\tfrac12F^2
&=\tfrac12(\mathbf E^2-\mathbf B^2)+i\,\mathbf E\cdot\mathbf B,\\
\tfrac12F F^{\mathrm R}
&=\tfrac12(\mathbf E^2+\mathbf B^2)+(\mathbf E\times\mathbf B)\cdot\sigv.
\end{align*}
These have direct standard physical interpretations: \(\tfrac12F^2\) packages the Lorentz invariants, while \(\tfrac12F F^{\mathrm R}\) packages the energy density and the Poynting vector. A simple bra-ket built from \(\rev{F}F\) is not the standard Maxwell-field pairing used here; the physically relevant quantities are the field bilinears together with the field equations and their gauge or subsidiary constraints.

The consistent summary is therefore this:
\begin{align*}
\text{spin}
&\neq
\text{component count alone},\\
\text{spin}
&=
(\text{carrier},\ \text{operator law},\ \text{physical bilinears},\ \text{constraints}).
\end{align*}
Within that summary, \((\cdot)^{\mathrm R}\) is the pairing used for quantum expectations, \((\cdot)^{*\mathrm R}\) is the metric pairing used for algebraic factorization, and the pure sigma-conjugate pairing remains an auxiliary companion unless and until a further physical identification is supplied.

\subsubsection{Spin-0 Wave Mechanics}
The scalar relativistic wave equation can be read directly from the mass-shell identity developed earlier in the kinematic section. Here \(\psi\in\C\) is a complex scalar field, \(m\in\R\), and the free-particle operators are defined using the differential conventions from Section \ref{sec:differential}. Start from
\begin{equation*}
\langle\!\langle mU\rangle\!\rangle^2=m^2
\qquad\Rightarrow\qquad
-E^2+\mathbf P^2+m^2=0.
\end{equation*}
For a free particle,
\begin{equation*}
E=i\partial_t,
\qquad
\mathbf P=-i\gradv,
\end{equation*}
and equivalently
\begin{equation*}
E+\mathbf P\cdot\sigv=i\partial^*.
\end{equation*}
Acting on a scalar field \(\psi\), the operator chain is
\begin{align*}
(-E^2+\mathbf P^2+m^2)\psi
&=\bigl(-(i\partial_t)^2-\gradv^2+m^2\bigr)\psi\notag\\
&=\bigl(\partial_t^2-\gradv^2+m^2\bigr)\psi\notag\\
&=(\partial\partial^*+m^2)\psi\notag\\
&=(\Box+m^2)\psi.
\end{align*}
Thus the free spin-0 relativistic wave equation is
\begin{equation*}
(\Box+m^2)\psi=0
\qquad\text{Klein--Gordon equation}.
\end{equation*}
This is the scalar wave equation recorded in Eq.\ \eqref{eq:kg-free}. The dispersion relation \(\omega^2=\mathbf k^2+m^2\) and the later positive/negative-frequency split are consequences of this same operator identity.

\subsubsection{Positive and Negative Frequency Sectors}
The frequency decomposition is one of the central interpretive steps in the scalar theory. Here \(\psi,\psi(\pm),c(\pm)\in\C\), while \(m\in\R\) and \(\gradv\) is the spatial derivative operator fixed earlier in Section \ref{sec:differential}. The frequency operator is
\begin{align*}
\Omega
&=\sqrt{m^2-\gradv^2}\notag\\
&=\langle\!\langle m+\sigv\cdot\gradv\rangle\!\rangle.
\end{align*}
Accordingly,
\begin{align*}
\Box+m^2
&=\partial_t^2+\Omega^2\notag\\
&=\partial_t^2+(m^2-\gradv^2)\notag\\
&=(\partial_t-i\Omega)(\partial_t+i\Omega)\notag\\
&=(\partial_t+i\Omega)(\partial_t-i\Omega).
\end{align*}
The positive- and negative-frequency sectors are then defined by
\begin{align*}
(\partial_t\pm i\Omega)\psi(\pm)&=0,\\
\psi(\pm)&=c(\pm)\exp(\mp i\Omega t),\\
\psi&=\psi(+)+\psi(-).
\end{align*}
This decomposition is not merely notation: it separates the two energy branches of the scalar theory in the same way that the Fourier transform separates positive and negative frequencies \cite{greiner2000}.

\subsubsection{Interpreting Root Operators with Fourier}
Square-root operators are easiest to interpret one Fourier mode at a time. In this subsection \(\psi\in\C\), the coefficients \(c_n\in\C\), the frequencies \(\omega_n\in\R\), and the wave vectors \(\mathbf k_n\in\R^3\). Write
\begin{align*}
\psi&=\sum_{n} c_{n}f_{n},\\
f_{n}&=\exp\!\bigl(i(\omega_{n}t+\mathbf k_{n}\cdot\mathbf r)\bigr).
\end{align*}
Then the d'Alembertian and its square root act modewise:
\begin{align*}
\Box\psi
&=\sum_{n}(\mathbf k_{n}^{2}-\omega_{n}^{2})c_{n}f_{n},\\
\pm\sqrt{\Box}\,\psi
&=\sum_{n}\pm\sqrt{\mathbf k_{n}^{2}-\omega_{n}^{2}}\,c_{n}f_{n},\\
(\Box+m^{2})\psi
&=\sum_{n}(\mathbf k_{n}^{2}+m^{2}-\omega_{n}^{2})c_{n}f_{n}.
\end{align*}
For on-shell Klein--Gordon modes,
\begin{equation*}
\mathbf k_{n}^{2}+m^{2}=\omega_{n}^{2},
\qquad
\omega_{n}=\pm\sqrt{\mathbf k_{n}^{2}+m^{2}},
\end{equation*}
so the operator factorization can be read as the modewise identity
\begin{equation*}
\Box+m^{2}=(m\mp i\sqrt{\Box})(m\pm i\sqrt{\Box}).
\end{equation*}
Likewise, if \(\Omega=\sqrt{m^{2}-\gradv^{2}}\) as in Eq.\ \eqref{eq:frequency-split}, then
\begin{align*}
\Omega^{2}\psi
&=\sum_{n}(m^{2}+\mathbf k_{n}^{2})c_{n}f_{n},\\
\Omega\psi
&=\sum_{n}\sqrt{m^{2}+\mathbf k_{n}^{2}}\,c_{n}f_{n}.
\end{align*}
This is the concrete meaning of the root operators used above: they do not represent a formal symbolic manipulation, but multiplication of each Fourier mode by the corresponding dispersion function.

\subsubsection{Klein--Gordon with Potential Fields}
To keep the minimally coupled scalar theory in one self-contained place, start from the charged-particle Lagrangian
\begin{equation*}
\mathcal L=-\frac{m}{\gamma}-q(V-\mathbf v\cdot\mathbf A).
\end{equation*}
Here \(E\) denotes the energy operator, not the electric field \(\mathbf E\). The Hamiltonian and canonical momentum are then
\begin{equation*}
i\partial_t=\mathcal H=E+qV,
\qquad
-i\gradv=\delv\mathcal L=\mathbf P+q\mathbf A,
\end{equation*}
so the covariant energy and momentum operators are
\begin{equation*}
E=i\partial_t-qV,
\qquad
\mathbf P=-i\gradv-q\mathbf A.
\end{equation*}

Acting on a smooth test field \(\psi\), the Leibniz rules needed below are
\begin{align*}
\partial_t(V\psi)&=(\partial_tV)\psi+V\partial_t\psi,\\
\partial_m(A_n\psi)&=(\partial_mA_n)\psi+A_n\partial_m\psi.
\end{align*}
From these one obtains
\begin{align*}
E^2\psi
&=E(i\partial_t\psi-qV\psi)\notag\\
&=\left(-\partial_t^2-iq\,\partial_tV-2iqV\partial_t+q^2V^2\right)\psi,\\
\mathbf P^2\psi
&=\mathbf P\cdot(-i\gradv\psi-q\mathbf A\psi)\notag\\
&=\left(-\gradv^2+iq\,\gradv\cdot\mathbf A+2iq\,\mathbf A\cdot\gradv+q^2\mathbf A^2\right)\psi.
\end{align*}

Recall also
\begin{align*}
\Box&=\partial_t^2-\gradv^2,\\
\Phi&=V+\mathbf A\cdot\sigv,\\
S&=\partial_tV+\gradv\cdot\mathbf A,\\
\langle\!\langle\partial\rangle\!\rangle^2&=\Box,\\
\langle\!\langle\Phi\rangle\!\rangle^2&=V^2-\mathbf A^2.
\end{align*}
Substituting the operator expansions gives the minimally coupled Klein--Gordon equation in explicit form:
\begin{align*}
0&=(-E^2+\mathbf P^2+m^2)\psi\notag\\
&=\Bigl(\Box+iq(\partial_tV+\gradv\cdot\mathbf A)+2iq(V\partial_t+\mathbf A\cdot\gradv)\notag\\
&\qquad\qquad-q^2(V^2-\mathbf A^2)+m^2\Bigr)\psi\notag\\
&=\Bigl(\Box+iqS-q^2\langle\!\langle\Phi\rangle\!\rangle^2+m^2+2iq(V\partial_t+\mathbf A\cdot\gradv)\Bigr)\psi.
\end{align*}

The electromagnetic fields and commutators enter through
\begin{align*}
\mathbf E&=-\partial_t\mathbf A-\gradv V,\\
\mathbf B&=\gradv\times\mathbf A,\\
F&=(\mathbf E+i\mathbf B)\cdot\sigv,
\end{align*}
and
\begin{align*}
[E,P_m]\psi
&=(i\partial_t-qV)(-i\partial_m\psi-qA_m\psi)\notag\\
&\quad-(-i\partial_m-qA_m)(i\partial_t\psi-qV\psi)\notag\\
&=-iq\bigl((\partial_tA_m)\psi+A_m\partial_t\psi\bigr)\notag\\
&\quad-iq\bigl((\partial_mV)\psi+V\partial_m\psi\bigr)+iqV\partial_m\psi+iqA_m\partial_t\psi\notag\\
&=-iq(\partial_tA_m+\partial_mV)\psi\notag\\
&=iqE_m\psi,\\
[P_m,P_n]\psi
&=(-i\partial_m-qA_m)\bigl(-i\partial_n\psi-qA_n\psi\bigr)\notag\\
&\quad-(-i\partial_n-qA_n)\bigl(-i\partial_m\psi-qA_m\psi\bigr)\notag\\
&=iq(\partial_mA_n-\partial_nA_m)\psi.
\end{align*}
Hence
\begin{equation*}
[E,\mathbf P]=iq\mathbf E,
\qquad
\mathbf P\times\mathbf P=iq\mathbf B.
\end{equation*}
Substituting those commutators into the ordered products yields
\begin{align*}
(E+\sigv\cdot\mathbf P)(E-\sigv\cdot\mathbf P)
&=E^2-\mathbf P^2\notag\\
&\qquad+\sigv\cdot\bigl(-[E,\mathbf P]-i\,\mathbf P\times\mathbf P\bigr)\notag\\
&=E^2-\mathbf P^2-iqF,\\
(E-\sigv\cdot\mathbf P)(E+\sigv\cdot\mathbf P)
&=E^2-\mathbf P^2\notag\\
&\qquad+\sigv\cdot\bigl([E,\mathbf P]-i\,\mathbf P\times\mathbf P\bigr)\notag\\
&=E^2-\mathbf P^2-iqF^*.
\end{align*}
Since
\begin{equation*}
E+\sigv\cdot\mathbf P=i\partial^*-q\Phi,
\qquad
E-\sigv\cdot\mathbf P=i\partial-q\Phi^*
=(E+\sigv\cdot\mathbf P)^{*\mathrm R},
\end{equation*}
their scalar paravector metric is
\begin{equation*}
\langle\!\langle E\pm\sigv\cdot\mathbf P\rangle\!\rangle^2
=E^2-\mathbf P^2+\sigv\cdot(\pm [E,\mathbf P]-i\,\mathbf P\times\mathbf P).
\end{equation*}
This is not an ordinary scalar operator: operator ordering still matters, so the two first-order products remain distinct.
Consequently,
\begin{align*}
-(E+\sigv\cdot\mathbf P)(E+\sigv\cdot\mathbf P)^{*\mathrm R}
&=(\partial^*+iq\Phi)(\partial+iq\Phi^*),\\
-(E+\sigv\cdot\mathbf P)^{*\mathrm R}(E+\sigv\cdot\mathbf P)
&=(\partial+iq\Phi^*)(\partial^*+iq\Phi).
\end{align*}
The same factorizations can be rewritten directly in derivative form:
\begin{align*}
(E+\sigv\cdot\mathbf P)(E-\sigv\cdot\mathbf P)\psi
&=-(\partial^*+iq\Phi)(\partial+iq\Phi^*)\psi\notag\\
&=-\Box\psi-iq\,\partial^*(\Phi^*\psi)\notag\\
&\qquad-iq\,\Phi\,\partial\psi\notag\\
&\qquad+q^2\langle\!\langle\Phi\rangle\!\rangle^2\psi\notag\\
&=-\Box\psi-iqS\psi+q^2\langle\!\langle\Phi\rangle\!\rangle^2\psi\notag\\
&\qquad-2iq(V\partial_t+\mathbf A\cdot\gradv)\psi-iqF\psi,\\
(E-\sigv\cdot\mathbf P)(E+\sigv\cdot\mathbf P)\psi
&=-(\partial+iq\Phi^*)(\partial^*+iq\Phi)\psi\notag\\
&=-\Box\psi-iq\,\partial(\Phi\psi)\notag\\
&\qquad-iq\,\Phi^*\,\partial^*\psi\notag\\
&\qquad+q^2\langle\!\langle\Phi\rangle\!\rangle^2\psi\notag\\
&=-\Box\psi-iqS\psi+q^2\langle\!\langle\Phi\rangle\!\rangle^2\psi\notag\\
&\qquad-2iq(V\partial_t+\mathbf A\cdot\gradv)\psi-iqF^*\psi.
\end{align*}
Equivalently,
\begin{align*}
\partial^*(\Phi^*\psi)+\Phi\,\partial\psi
&=S\psi+2(V\partial_t+\mathbf A\cdot\gradv)\psi+F\psi,\\
\partial(\Phi\psi)+\Phi^*\,\partial^*\psi
&=S\psi+2(V\partial_t+\mathbf A\cdot\gradv)\psi+F^*\psi,
\end{align*}
and their difference isolates the electric contribution:
\begin{align*}
\bigl(\partial^*(\Phi^*\psi)+\Phi\,\partial\psi\bigr)
&-\bigl(\partial(\Phi\psi)+\Phi^*\,\partial^*\psi\bigr)\notag\\
&=\bigl([\Phi,\partial]+[\partial,\Phi]^*\bigr)\psi\notag\\
&=(F-F^*)\psi\notag\\
&=2(\mathbf E\cdot\sigv)\psi.
\end{align*}

In summary,
\begin{align*}
E&:=i\partial_t-qV, &
\mathbf P&:=-i\gradv-q\mathbf A,\\
\partial&:=\partial_t+\sigv\cdot\gradv, &
\partial^*&:=\partial_t-\sigv\cdot\gradv,\\
\Phi&:=V+\mathbf A\cdot\sigv, &
F&:=(\mathbf E+i\mathbf B)\cdot\sigv,\\
S&:=\partial_tV+\gradv\cdot\mathbf A, &
\mathbf E&:=-\partial_t\mathbf A-\gradv V,\\
\mathbf B&:=\gradv\times\mathbf A,
\end{align*}
and the scalar Klein--Gordon operator can be recovered from either ordered first-order product:
\begin{align*}
0&=(-E^2+\mathbf P^2+m^2)\psi\notag\\
&=\Bigl(\Box+iqS-q^2\langle\!\langle\Phi\rangle\!\rangle^2+m^2+2iq(V\partial_t+\mathbf A\cdot\gradv)\Bigr)\psi\notag\\
&=(\partial^*+iq\Phi)(\partial+iq\Phi^*)\psi-iqF\psi+m^2\psi\notag\\
&=(\partial+iq\Phi^*)(\partial^*+iq\Phi)\psi-iqF^*\psi+m^2\psi.
\end{align*}
The ordered first-order products themselves obey
\begin{align*}
(E\pm\sigv\cdot\mathbf P)(E\mp\sigv\cdot\mathbf P)
&=E^2-(\sigv\cdot\mathbf P)^2\mp[E,\sigv\cdot\mathbf P]\notag\\
&=E^2-\mathbf P^2-iq(\pm\mathbf E+i\mathbf B)\cdot\sigv.
\end{align*}
Taking the scalar part therefore recovers
\begin{equation*}
0=\eqtext{scalar}\!\bigl((\partial+iq\Phi^*)(\partial^*+iq\Phi)\bigr)\psi+m^2\psi.
\end{equation*}
Only the scalar part reproduces the Klein--Gordon operator itself; the full ordered products retain the non-scalar field terms \(-iqF\) and \(-iqF^*\) that distinguish the two factorizations.

\subsubsection{Spin-\texorpdfstring{$\frac12$}{1/2} Wave Mechanics}
Before the electromagnetic field is turned on, the same auxiliary \(2\times2\) map from Eq.\ \eqref{eq:W-map} gives a self-contained first-order version of the relativistic mass shell. Consider the two-component field
\begin{equation*}
\bPsi=
\begin{Bmatrix}
\bPsi_1\\[2pt]
\bPsi_2
\end{Bmatrix}.
\end{equation*}
For any complex paravector or paravector-valued operator \(Z\), define
\begin{align*}
\Wmat(Z)&:=
\begin{bmatrix}
0 & Z \\
\crev{Z} & 0
\end{bmatrix},\\
\Wmat(Z)^2&=\operatorname{diag}(Z\crev{Z},\crev{Z}Z).
\end{align*}
If \(a,b\in\R\), then \(\Wmat((a+ib)Z)=(a+ib)\Wmat(Z)\). In particular,
\begin{align*}
\Wmat(i\partial)
&=i
\begin{bmatrix}
0 & \partial \\
\partial^* & 0
\end{bmatrix}
=i\Wmat(\partial),\\
\Wmat(i\partial)^2
&=-\operatorname{diag}(\partial\partial^*,\partial^*\partial)\notag\\
&=-\Box\I.
\end{align*}

Define the free Dirac operator matrix and its mass conjugate by
\begin{equation*}
\Dmat:=\Wmat(i\partial)-m\I,
\qquad
\Dmat':=\Wmat(i\partial)+m\I.
\end{equation*}
Then
\begin{align*}
\Dmat'\Dmat
&=(\Wmat(i\partial)+m\I)(\Wmat(i\partial)-m\I)\notag\\
&=\Wmat(i\partial)^2-m^2\I\notag\\
&=-(\Box+m^2)\I.
\end{align*}
Thus the free Dirac equation
\begin{equation*}
\Dmat\bPsi=0
\end{equation*}
forces each component to satisfy the Klein--Gordon operator. Formally, its characteristic polynomial is
\begin{align*}
0
&=\det(\Dmat-s\I)\notag\\
&=\det
\begin{bmatrix}
-(s+m) & i\partial \\
i\partial^* & -(s+m)
\end{bmatrix}\notag\\
&=(s+m)^2+\Box,
\end{align*}
so the modewise eigenvalue branches are \(s=-m\pm i\sqrt{\Box}\), with the square-root operator interpreted in the Fourier sense from the previous subsection.

Writing out \(\Dmat\bPsi=0\) gives the two coupled equations
\begin{align*}
i\partial\bPsi_2-m\bPsi_1&=0,\\
i\partial^*\bPsi_1-m\bPsi_2&=0.
\end{align*}
Solving the first equation for the upper component and substituting into the second yields
\begin{align*}
\bPsi_1&=\frac{i\partial\bPsi_2}{m},\\
0&=i\partial^*\bPsi_1-m\bPsi_2\notag\\
&=-\frac{\partial^*\partial\bPsi_2}{m}-m\bPsi_2\notag\\
&=-\frac{(\Box+m^2)\bPsi_2}{m}.
\end{align*}
Hence the lower component is a Klein--Gordon solution, and the upper component is obtained from it by one first-order derivative.

Conversely, let \(\psi'\) be any scalar Klein--Gordon solution and let \(a,b\in\C\) be constants. Define
\begin{equation*}
\bPsi_1=\frac{(ma+ib\,\partial)\psi'}{m},
\qquad
\bPsi_2=\frac{(ia\,\partial^*+mb)\psi'}{m}.
\end{equation*}
Because \(a\) and \(b\) are constant and \((\Box+m^2)\psi'=0\),
\begin{align*}
i\partial\bPsi_2-m\bPsi_1
&=\frac{(-\partial a\,\partial^*+im\,\partial b-m^2a-imb\,\partial)\psi'}{m}\notag\\
&=-\frac{a}{m}(\Box+m^2)\psi'=0,\\
i\partial^*\bPsi_1-m\bPsi_2
&=\frac{(mi\,\partial^*a-\partial^*b\,\partial-mi\,a\,\partial^*-m^2b)\psi'}{m}\notag\\
&=-\frac{b}{m}(\Box+m^2)\psi'=0.
\end{align*}
The family generated from a scalar Klein--Gordon solution therefore solves the free two-component system for any constant \(a\) and \(b\). The simplest choice is \(a=1\) and \(b=0\), for which
\begin{equation*}
\bPsi=
\begin{Bmatrix}
\psi'\\[2pt]
\dfrac{i}{m}\,\partial^*\psi'
\end{Bmatrix}.
\end{equation*}
This is the free spin-\(\tfrac12\) wave-mechanics factorization in its most explicit form: the Dirac operator linearizes the Klein--Gordon mass shell, and scalar solutions generate two-component spinor solutions by a first-order lift.

\subsubsection{Dirac with Potential Field}
Consider again the two-component field
\begin{equation*}
\bPsi=
\begin{Bmatrix}
\bPsi_1\\[2pt]
\bPsi_2
\end{Bmatrix}.
\end{equation*}
For any complex paravector or paravector-valued operator \(Z\), define
\begin{align*}
\Wmat(Z)&:=
\begin{bmatrix}
0 & Z \\
\crev{Z} & 0
\end{bmatrix},\\
\Wmat(Z)^2&=\operatorname{diag}(Z\crev{Z},\crev{Z}Z),\\
\Wmat(Z)^*&=\Wmat(Z^*),\\
\Wmat(Z)^{\mathrm R}&=\Wmat(Z^{\mathrm R}).
\end{align*}
The Dirac operator with potential is
\begin{align*}
\Dmat&:=\Wmat(i\partial-q\Phi^*)-m\I\notag\\
&=i\Wmat(\partial)-q\Wmat(\Phi^*)-m\I.
\end{align*}
Because
\begin{align*}
\Wmat(\partial)^2&=\Box\I,\\
\Wmat(\partial)\Wmat(\Phi^*)
&=
\begin{bmatrix}
0 & \partial\\
\partial^* & 0
\end{bmatrix}
\begin{bmatrix}
0 & \Phi^*\\
\Phi & 0
\end{bmatrix}
=\operatorname{diag}(\partial\Phi,\partial^*\Phi^*),\\
\Wmat(\Phi^*)\Wmat(\partial)
&=
\begin{bmatrix}
0 & \Phi^*\\
\Phi & 0
\end{bmatrix}
\begin{bmatrix}
0 & \partial\\
\partial^* & 0
\end{bmatrix}
=\operatorname{diag}(\Phi^*\partial^*,\Phi\partial),\\
\Wmat(\Phi^*)^2
&=
\begin{bmatrix}
0 & \Phi^*\\
\Phi & 0
\end{bmatrix}^{2}
=\operatorname{diag}(\Phi^*\Phi,\Phi\Phi^*)\notag\\
&=\langle\!\langle\Phi\rangle\!\rangle^2\I,
\end{align*}
one obtains
\begin{align*}
\Wmat(i\partial-q\Phi^*)^2
&=-\Wmat(\partial)^2\notag\\
&\qquad-iq\bigl(\Wmat(\partial)\Wmat(\Phi^*)+\Wmat(\Phi^*)\Wmat(\partial)\bigr)\notag\\
&\qquad+q^2\Wmat(\Phi^*)^2\notag\\
&=-\Box\I-iq\,\operatorname{diag}(\partial\Phi+\Phi^*\partial^*,\partial^*\Phi^*+\Phi\partial)\notag\\
&\qquad+q^2\langle\!\langle\Phi\rangle\!\rangle^2\I\notag\\
&=-\operatorname{diag}\!\Big(
(\partial+iq\Phi^*)(\partial^*+iq\Phi),\notag\\
&\qquad\qquad\qquad\quad
(\partial^*+iq\Phi)(\partial+iq\Phi^*)
\Big).
\end{align*}
The Dirac equation
\begin{equation*}
\Dmat\bPsi=0
\end{equation*}
therefore becomes the two coupled equations
\begin{align*}
-m\bPsi_1+(i\partial-q\Phi^*)\bPsi_2&=0,\\
(i\partial^*-q\Phi)\bPsi_1-m\bPsi_2&=0.
\end{align*}
Solving the first equation for the upper component and substituting into the second gives
\begin{align*}
\bPsi_1&=\frac{(i\partial-q\Phi^*)\bPsi_2}{m},\\
0&=(i\partial^*-q\Phi)\bPsi_1-m\bPsi_2\notag\\
&=-\frac{\bigl((\partial^*+iq\Phi)(\partial+iq\Phi^*)+m^2\bigr)\bPsi_2}{m}.
\end{align*}
Thus each component satisfies a Klein--Gordon-type operator with the appropriate ordered field coupling.

The mass-conjugate operator is
\begin{equation*}
\Dmat':=\Wmat(i\partial-q\Phi^*)+m\I.
\end{equation*}
Multiplying by \(\Dmat'\) yields
\begin{align*}
0=\Dmat'\Dmat\bPsi
&=\bigl(\Wmat(i\partial-q\Phi^*)^2-m^2\I\bigr)\bPsi\notag\\
&=-\operatorname{diag}\!\Big(
(\partial+iq\Phi^*)(\partial^*+iq\Phi)+m^2,\notag\\
&\qquad\qquad\qquad\quad
(\partial^*+iq\Phi)(\partial+iq\Phi^*)+m^2
\Big)\bPsi.
\end{align*}
Equivalently,
\begin{align*}
0&=\bigl((\partial+iq\Phi^*)(\partial^*+iq\Phi)+m^2\bigr)\bPsi_1\notag\\
&=\bigl(-E^2+\mathbf P^2+m^2+iqF^*\bigr)\bPsi_1,\\
0&=\bigl((\partial^*+iq\Phi)(\partial+iq\Phi^*)+m^2\bigr)\bPsi_2\notag\\
&=\bigl(-E^2+\mathbf P^2+m^2+iqF\bigr)\bPsi_2.
\end{align*}
Once the square is taken, the coupled Dirac system separates into two Klein--Gordon-type equations whose field terms differ only by the conjugation pattern fixed by the operator ordering.

\subsubsection{Gauge Transform of the Dirac System}
Let \(\lambda\in\R\) be a real scalar gauge field and define
\begin{equation*}
\Phi'=\Phi+\partial^*\lambda.
\end{equation*}
Because \(\lambda\) is real,
\begin{equation*}
\partial\lambda=(\Phi'-\Phi)^*.
\end{equation*}
Transform the two-component field by
\begin{equation*}
\bPsi=e^{iq\lambda}\bPsi'.
\end{equation*}
Then the derivative of the phase factor is
\begin{align*}
\partial\bigl(e^{iq\lambda}\bPsi'\bigr)
&=e^{iq\lambda}\bigl(iq\,\partial\lambda+\partial\bigr)\bPsi'.
\end{align*}
Therefore,
\begin{align*}
0&=(\Wmat(i\partial-q\Phi^*)-m\I)\bPsi\notag\\
&=(\Wmat(i\partial-q\Phi^*)-m\I)e^{iq\lambda}\bPsi'\notag\\
&=e^{iq\lambda}(\Wmat(i\partial-q\,\partial\lambda-q\Phi^*)-m\I)\bPsi'\notag\\
&=e^{iq\lambda}(\Wmat(i\partial-q(\Phi'-\Phi)^*-q\Phi^*)-m\I)\bPsi'\notag\\
&=e^{iq\lambda}(\Wmat(i\partial-q\Phi'^*)-m\I)\bPsi'.
\end{align*}
Hence
\begin{equation*}
(\Wmat(i\partial-q\Phi'^*)-m\I)\bPsi'=0.
\end{equation*}
This is the gauge-covariance statement: the phase change of the Dirac field is exactly compensated by the shift of the potential, so the coupled Dirac equation keeps the same form in the primed gauge.

\subsubsection{Dirac Equation Components}
For the component form of the coupled Dirac equation, write
\begin{align*}
\partial&:=\partial_t+\sigv\cdot\gradv,\\
\partial^*&:=\partial_t-\sigv\cdot\gradv,\\
\Phi&:=V+\mathbf A\cdot\sigv\in\R\oplus\R^3,\\
\Phi^*&:=V-\mathbf A\cdot\sigv,\\
\psi&:=\psi_s+\psi_{\mathbf v}\cdot\sigv\in\C\oplus\C^3.
\end{align*}
For the field term, use
\begin{align*}
F&:=\mathbf F\cdot\sigv,\\
\eqtext{scalar}(F)&=0,\\
\mathbf F&:=\mathbf E+i\mathbf B\in\C^3,\\
F^*&=-\mathbf F^*\cdot\sigv.
\end{align*}
The supporting tables then give the scalar and vector parts of \(\partial\psi\), \(\partial^*\psi\), \(\Phi\psi\), \(\Phi^*\psi\), \(iqF\psi\), and \(iqF^*\psi\) explicitly.
\noindent The direct component action of \(\pv{\partial}\) and \(\conj{\pv{\partial}}\) on a complex paravector test field is recorded in Tables \ref{tab:components-partial-psi} and \ref{tab:components-partialstar-psi}.
\begin{componenttable}[tab:components-partial-psi]{\pv{\partial} \pv{\psi}_{\eqtext{block}}}
\componentrow{scalar}{\partial_t\eqfunc{scalar}(\pv{\psi}_{\eqtext{block}})+\gradv\cdot\eqfunc{vector}(\pv{\psi}_{\eqtext{block}})}{}
\componentrow{vector}{\partial_t\eqfunc{vector}(\pv{\psi}_{\eqtext{block}})+\gradv\eqfunc{scalar}(\pv{\psi}_{\eqtext{block}})+i(\gradv\times\eqfunc{vector}(\pv{\psi}_{\eqtext{block}}))}{}
\end{componenttable}

\begin{componenttable}[tab:components-partialstar-psi]{\conj{\pv{\partial}}\pv{\psi}_{\eqtext{block}}}
\componentrow{scalar}{\partial_t\eqfunc{scalar}(\pv{\psi}_{\eqtext{block}})-\gradv\cdot\eqfunc{vector}(\pv{\psi}_{\eqtext{block}})}{}
\componentrow{vector}{\partial_t\eqfunc{vector}(\pv{\psi}_{\eqtext{block}})-\gradv\eqfunc{scalar}(\pv{\psi}_{\eqtext{block}})-i(\gradv\times\eqfunc{vector}(\pv{\psi}_{\eqtext{block}}))}{}
\end{componenttable}

\noindent The component actions of the potential \(\pv{\Phi}\) and its conjugate on the same field are then shown in Tables \ref{tab:components-phi-psi} and \ref{tab:components-phistar-psi}.
\begin{componenttable}[tab:components-phi-psi]{\pv{\Phi}\pv{\psi}_{\eqtext{block}}}
\componentrow{scalar}{V\eqfunc{scalar}(\pv{\psi}_{\eqtext{block}})+\vct{A}\cdot\eqfunc{vector}(\pv{\psi}_{\eqtext{block}})}{}
\componentrow{vector}{V\eqfunc{vector}(\pv{\psi}_{\eqtext{block}})+\eqfunc{scalar}(\pv{\psi}_{\eqtext{block}})\,\vct{A}+i(\vct{A}\times\eqfunc{vector}(\pv{\psi}_{\eqtext{block}}))}{}
\end{componenttable}

\begin{componenttable}[tab:components-phistar-psi]{\conj{\pv{\Phi}}\pv{\psi}_{\eqtext{block}}}
\componentrow{scalar}{V\eqfunc{scalar}(\pv{\psi}_{\eqtext{block}})-\vct{A}\cdot\eqfunc{vector}(\pv{\psi}_{\eqtext{block}})}{}
\componentrow{vector}{V\eqfunc{vector}(\pv{\psi}_{\eqtext{block}})-\eqfunc{scalar}(\pv{\psi}_{\eqtext{block}})\,\vct{A}-i(\vct{A}\times\eqfunc{vector}(\pv{\psi}_{\eqtext{block}}))}{}
\end{componenttable}

\noindent The Dirac-block operator tables are completed by Tables
\ref{tab:components-iqf-psi} and \ref{tab:components-iqfstar-psi}, in which
the field-coupling terms \(iq\pv{F}\pv{\psi}_{\eqtext{block}}\) and
\(iq\conj{\pv{F}}\pv{\psi}_{\eqtext{block}}\) are recorded. In the second table,
\(\conj{\pv{F}}=-\vct{F}^{*}\cdot\sigv
=-(\vct{E}-i\vct{B})\cdot\sigv\), so no additional conjugation operation is
hidden in the notation.
\begin{componenttable}[tab:components-iqf-psi]{iq\pv{F}\pv{\psi}_{\eqtext{block}}}
\componentrow{scalar}{iq\bigl(\vct{F}\cdot\eqfunc{vector}(\pv{\psi}_{\eqtext{block}})\bigr)}{}
\componentrow{vector}{iq\,\eqfunc{scalar}(\pv{\psi}_{\eqtext{block}})\,\vct{F}
-q\bigl(\vct{F}\times\eqfunc{vector}(\pv{\psi}_{\eqtext{block}})\bigr)}{}
\end{componenttable}

\begin{componenttable}[tab:components-iqfstar-psi]{iq\conj{\pv{F}}\pv{\psi}_{\eqtext{block}}}
\componentrow{scalar}{-iq\bigl(\vct{F}^{*}\cdot\eqfunc{vector}(\pv{\psi}_{\eqtext{block}})\bigr)}{}
\componentrow{vector}{-iq\,\eqfunc{scalar}(\pv{\psi}_{\eqtext{block}})\,\vct{F}^{*}
+q\bigl(\vct{F}^{*}\times\eqfunc{vector}(\pv{\psi}_{\eqtext{block}})\bigr)}{}
\end{componenttable}


\subsubsection{Low-Energy Spin-0 Limit}
For the minimally coupled scalar theory, use the covariant operators
\begin{equation*}
E=i\partial_t-qV,
\qquad
\mathbf P=-i\gradv-q\mathbf A.
\end{equation*}
Remove the fast rest-mass oscillation by writing
\begin{equation*}
\psi'=e^{-imt}\psi.
\end{equation*}
Then
\begin{align*}
E^2\psi'&=e^{-imt}(E+m)^2\psi,\\
\mathbf P^2\psi'&=e^{-imt}\mathbf P^2\psi.
\end{align*}
Substituting into the minimally coupled Klein--Gordon equation gives
\begin{align*}
0&=(-E^2+\mathbf P^2+m^2)\psi'\notag\\
&=e^{-imt}\bigl(-(E+m)^2+\mathbf P^2+m^2\bigr)\psi,
\end{align*}
hence
\begin{equation*}
0=(-E^2-2mE+\mathbf P^2)\psi.
\end{equation*}
Formally iterating this relation gives
\begin{align*}
E&=\frac{\mathbf P^2-E^2}{2m},\\
&=\frac{\mathbf P^2}{2m}-\frac{\mathbf P^4}{8m^3}+\cdots.
\end{align*}
At low energy,
\begin{equation*}
E\psi\approx \frac{\mathbf P^2}{2m}\psi.
\end{equation*}
Solving for \(i\partial_t\psi\) yields the Schr\"odinger equation
\begin{equation*}
i\partial_t\psi
=\left(\frac{(i\gradv+q\mathbf A)^2}{2m}+qV\right)\psi.
\end{equation*}
If \(\mathbf A=0\), this reduces to
\begin{equation*}
i\partial_t\psi
=\left(-\frac{\gradv^2}{2m}+qV\right)\psi.
\end{equation*}
The scalar low-energy limit is therefore already encoded in the relativistic Klein--Gordon theory.

\subsubsection{Low-Energy Spin-\texorpdfstring{$\frac12$}{1/2} and Pauli Limit}
Start from the two-component field
\begin{equation*}
\bPsi'=
\begin{Bmatrix}
\bPsi_1'\\[2pt]
\bPsi_2'
\end{Bmatrix}
\end{equation*}
and the potential-coupled Dirac equation
\begin{equation*}
\bigl(\Wmat(i\partial-q\Phi^*)-m\I\bigr)\bPsi'=0.
\end{equation*}
Its two component equations are
\begin{align*}
(i\partial-q\Phi^*)\bPsi_2'-m\bPsi_1'&=0,\\
(i\partial^*-q\Phi)\bPsi_1'-m\bPsi_2'&=0.
\end{align*}
Write
\begin{align*}
\partial&:=\partial_t+\sigv\cdot\gradv,
&
\partial^*&:=\partial_t-\sigv\cdot\gradv,\\
\Phi&:=V+\mathbf A\cdot\sigv,
&
\Phi^*&:=V-\mathbf A\cdot\sigv.
\end{align*}
Then the two equations become
\begin{align*}
(i\partial_t-qV)\bPsi_2'
+\bigl((i\gradv+q\mathbf A)\cdot\sigv\bigr)\bPsi_2'
-m\bPsi_1'&=0,\\
(i\partial_t-qV)\bPsi_1'
-\bigl((i\gradv+q\mathbf A)\cdot\sigv\bigr)\bPsi_1'
-m\bPsi_2'&=0.
\end{align*}
Introduce the covariant operators
\begin{equation*}
E:=i\partial_t-qV,
\qquad
\mathbf P:=-i\gradv-q\mathbf A.
\end{equation*}
The system is then
\begin{align*}
(E-\mathbf P\cdot\sigv)\bPsi_2'-m\bPsi_1'&=0,\\
(E+\mathbf P\cdot\sigv)\bPsi_1'-m\bPsi_2'&=0.
\end{align*}
Peel off the fast rest-mass oscillation with
\begin{equation*}
\bPsi_n'=e^{-imt}\bPsi_n.
\end{equation*}
Since
\begin{equation*}
E\bPsi_n'=e^{-imt}(E+m)\bPsi_n,
\end{equation*}
the reduced system becomes
\begin{align*}
(E+m-\mathbf P\cdot\sigv)\bPsi_2-m\bPsi_1&=0,\\
(E+m+\mathbf P\cdot\sigv)\bPsi_1-m\bPsi_2&=0.
\end{align*}
Solve the first equation for the upper component:
\begin{equation*}
\bPsi_1=\frac{(E+m-\mathbf P\cdot\sigv)\bPsi_2}{m}.
\end{equation*}
Substituting this into the second equation and multiplying by \(m\) gives
\begin{equation*}
0=\bigl((E+m+\mathbf P\cdot\sigv)(E+m-\mathbf P\cdot\sigv)-m^2\bigr)\bPsi_2.
\end{equation*}
Expand the operator product as
\begin{align*}
&(E+m+\mathbf P\cdot\sigv)\notag\\
&\qquad\times(E+m-\mathbf P\cdot\sigv)-m^2\notag\\
&=(E+m)^2-m^2-(\mathbf P\cdot\sigv)^2\notag\\
&\qquad+[\mathbf P\cdot\sigv,E].
\end{align*}
Now
\begin{align*}
-(\mathbf P\cdot\sigv)^2
&=-\mathbf P^2-i(\mathbf P\times\mathbf P)\cdot\sigv,\\
[P_m,P_n]
&=[-i\partial_m-qA_m,-i\partial_n-qA_n]\notag\\
&=iq(\partial_mA_n-\partial_nA_m),\\
\mathbf P\times\mathbf P
&:=([P_2,P_3],-[P_1,P_3],[P_1,P_2])\notag\\
&=iq\,\gradv\times\mathbf A\notag\\
&=iq\,\mathbf B,
\end{align*}
so
\begin{equation*}
-i(\mathbf P\times\mathbf P)\cdot\sigv=q\,\mathbf B\cdot\sigv.
\end{equation*}
Also,
\begin{align*}
[\mathbf P\cdot\sigv,E+m]
&=[\mathbf P,E]\cdot\sigv\notag\\
&=iq(\gradv V+\partial_t\mathbf A)\cdot\sigv\notag\\
&=-iq\,\mathbf E\cdot\sigv,
\end{align*}
where
\begin{equation*}
\mathbf E:=-\partial_t\mathbf A-\gradv V,
\qquad
\mathbf B:=\gradv\times\mathbf A.
\end{equation*}
Defining
\begin{equation*}
F:=(\mathbf E+i\mathbf B)\cdot\sigv,
\end{equation*}
one has
\begin{equation*}
-iqF=q\,\mathbf B\cdot\sigv-iq\,\mathbf E\cdot\sigv.
\end{equation*}
Therefore the exact lower-component equation is
\begin{align*}
0
&=\bigl((E+m)^2-m^2-\mathbf P^2+q\,\mathbf B\cdot\sigv+[\mathbf P\cdot\sigv,E]\bigr)\bPsi_2\notag\\
&=\bigl(E^2+2mE-\mathbf P^2+q\,\mathbf B\cdot\sigv-iq\,\mathbf E\cdot\sigv\bigr)\bPsi_2\notag\\
&=\bigl(E^2+2mE-\mathbf P^2-iqF\bigr)\bPsi_2.
\end{align*}
This can be iterated recursively as
\begin{align*}
E
&=\frac{\mathbf P^2-E^2-q\,\mathbf B\cdot\sigv-[\mathbf P\cdot\sigv,E]}{2m}\notag\\
&=\frac{\mathbf P^2-q\,\mathbf B\cdot\sigv}{2m}+O(1/m^2).
\end{align*}
At low energy, the dominant terms satisfy
\begin{equation*}
0\approx(2mE-\mathbf P^2+q\,\mathbf B\cdot\sigv)\bPsi_2.
\end{equation*}
Solving for \(i\partial_t\bPsi_2\) yields the Pauli equation
\begin{equation*}
i\partial_t\bPsi_2
=\left(\frac{(-i\gradv-q\mathbf A)^2}{2m}+qV\right)\bPsi_2
-\frac{q}{2m}\,(\mathbf B\cdot\sigv)\bPsi_2.
\end{equation*}
The nonrelativistic spin coupling is therefore not an extra postulate: it is the leading field correction that survives when the massive Dirac system is expanded around its rest-energy oscillation \cite{pauli1927,sakurai2017}.

\subsubsection{Spin Up and Spin Down}
Let the magnetic-field direction define the spin axis,
\begin{equation*}
\mathbf u:=\frac{\mathbf B}{\|\mathbf B\|},
\end{equation*}
and use the projection operators
\begin{align*}
\Pi(\pm)&=\tfrac12(1\pm \mathbf u\cdot\sigv),\\
\Pi(\pm)^2&=\Pi(\pm),\\
\Pi(+)\Pi(-)&=0,\\
\Pi(+)+\Pi(-)&=1,\\
(\mathbf u\cdot\sigv)\Pi(\pm)&=\pm\Pi(\pm),\\
\Pi(\pm)^{\mathrm R}&=\Pi(\pm),\\
\Pi(\pm)^*&=\Pi(\mp).
\end{align*}
Accordingly,
\begin{equation*}
\mathbf u\cdot\sigv=(+1)\Pi(+)+(-1)\Pi(-).
\end{equation*}
Decompose the Pauli state by
\begin{align*}
\bPsi(\pm)&:=\Pi(\pm)\bPsi,\\
(\mathbf u\cdot\sigv)\bPsi(\pm)&=\pm\bPsi(\pm),\\
\bPsi&=\bPsi(+)+\bPsi(-),
\end{align*}
so \(\bPsi(+)\) and \(\bPsi(-)\) are the spin-up and spin-down parts along the field axis.

Write the Pauli equation as
\begin{align*}
\mathbf P&:=-i\gradv-q\mathbf A,\\
i\partial_t\bPsi
&=\left(\frac{\mathbf P^2}{2m}+qV\right)\bPsi
-\frac{q}{2m}\,\|\mathbf B\|(\mathbf u\cdot\sigv)\bPsi.
\end{align*}
Equivalently,
\begin{equation*}
i\partial_t\bPsi=H\bPsi,
\end{equation*}
with
\begin{equation*}
H:=\frac{\mathbf P^2}{2m}+qV-\frac{q}{2m}\,\|\mathbf B\|\,\mathbf u\cdot\sigv.
\end{equation*}
Applying \(\Pi(\pm)\) to the time evolution gives
\begin{align*}
i\partial_t\bPsi(\pm)
&=i(\partial_t\Pi(\pm))\bPsi+\Pi(\pm)(i\partial_t\bPsi)\notag\\
&=i(\partial_t\Pi(\pm))\bPsi+\Pi(\pm)(H\bPsi).
\end{align*}
Add zero in the form \(H\Pi(\pm)\bPsi-H\Pi(\pm)\bPsi\):
\begin{align*}
\Pi(\pm)(H\bPsi)
&=\Pi(\pm)(H\bPsi)+H\Pi(\pm)\bPsi-H\Pi(\pm)\bPsi\notag\\
&=H\bPsi(\pm)+[\Pi(\pm),H]\bPsi.
\end{align*}
Therefore
\begin{align*}
i\partial_t\bPsi(\pm)
&=H\bPsi(\pm)+\bigl(i\partial_t\Pi(\pm)+[\Pi(\pm),H]\bigr)\bPsi\notag\\
&=:H\bPsi(\pm)+S(\pm).
\end{align*}
In explicit form,
\begin{align*}
i\partial_t\bPsi(\pm)
&=\left(\frac{\mathbf P^2}{2m}+qV\right)\bPsi(\pm)\notag\\
&\qquad-\frac{q}{2m}\,\|\mathbf B\|(\mathbf u\cdot\sigv)\bPsi(\pm)+S(\pm),
\end{align*}
with coupling term
\begin{equation*}
S(\pm):=\left(i\partial_t\Pi(\pm)+\frac{[\mathbf P^2,\Pi(\pm)]}{2m}\right)\bPsi.
\end{equation*}
If \(\mathbf B\) is uniform and static, then \(\Pi(\pm)\) is constant and
\begin{equation*}
S(+)=S(-)=0,
\end{equation*}
so the two spin sectors decouple. In that case
\begin{equation*}
(-\mathbf B\cdot\sigv)\bPsi(\pm)=\mp \|\mathbf B\|\,\bPsi(\pm),
\end{equation*}
and the Zeeman effect appears as
\begin{equation*}
i\partial_t\bPsi(\pm)
=\left(\frac{\mathbf P^2\mp q\|\mathbf B\|}{2m}+qV\right)\bPsi(\pm).
\end{equation*}
The field direction therefore selects a preferred axis, and the two projector parts evolve with opposite Zeeman shifts.
